\documentclass[11pt]{letter}

\usepackage[margin=0.5in]{geometry}
\usepackage{hyperref}

\begin{document}
\begin{letter}{}
\opening{Dear Hiring Team,}

I am writing to apply for the Machine Learning Engineer, New Grad position on Quora's content quality and safety team. I completed my M.Sc. in Computer Science at the University of Calgary in February 2026, and my background combines Python/PyTorch research, backend development for NLP services, and hands-on LLM workflow automation. I am particularly interested in the evaluation side of this role: understanding where an agent's decisions fail, testing changes carefully, and making the system more useful without losing sight of reliability or computational cost.

My graduate ML work involved developing a super-resolution framework for digital elevation data. I implemented and trained deep ResNet models in PyTorch, built data acquisition and preprocessing pipelines, and evaluated reconstruction quality using both numerical and terrain-specific metrics. One important lesson was that random data splits could make performance look better than it would be on unfamiliar regions. I used region-based validation to get a more realistic view of generalization. That experience shaped how I approach model evaluation: start with the conditions the system will actually face, then examine what an aggregate score may be hiding. I would bring that approach to evaluating agents that encounter changing content and abuse patterns.

I am also building an agentic job-search platform that uses LLM workflows for job discovery, resume matching, skill-gap analysis, and application decision support. I prototyped the workflow in n8n and have been moving toward a more flexible OpenClaw-based system, integrating Google Sheets, Google Drive, and Telegram. The project has given me practical experience connecting model-driven tasks with external tools and structured records. I also use Codex for AI-assisted development. This makes the opportunity to build and iterate on LLM-powered moderation workflows especially relevant to the work I am already pursuing.

My software engineering experience would help me contribute beyond model experimentation. At Asr Gooyesh Pardaz, I developed Django REST APIs backed by PostgreSQL for NLP and speech services, built authentication and payment-integration features, and contributed to a chatbot using GPT-style responses and FAQ data. In my graduate work with BigGeo and Mitacs, I built C++/Python data systems and improved encoding performance through multithreading, reducing processing time for ten tensors from 233 seconds to 26 seconds. These projects taught me to connect algorithmic choices with implementation, data handling, and measured system performance.

Teaching has also influenced how I work with a team. At the University of Calgary, I supported students working on data pipelines, databases, and large-scale processing. At Neuromatch Academy, I mentored an international group through PyTorch tutorials and deep-learning projects. Explaining methods and helping others debug has made me more deliberate about documenting assumptions and communicating technical decisions.

Quora's role appeals to me because the quality of automated decisions directly affects whether people can participate in a useful, welcoming space for sharing knowledge. I would welcome the chance to bring my ML evaluation, agent-workflow, and software development experience to that work, while learning from engineers operating these systems at Quora's scale.

Thank you for considering my application. I look forward to discussing how I could contribute to the team.

\closing{Sincerely,

Amir Mirzai Golpayegani

\href{mailto:amirgolpaa24@gmail.com}{amirgolpaa24@gmail.com}}
\end{letter}
\end{document}
